%!TEX engine = xelatex
\documentclass[11pt,a4paper]{article}
\usepackage{iftex}
\RequireXeTeX
\usepackage{fontspec}
\usepackage{polyglossia}
\usepackage{tikz}
\usepackage{xcolor}
\setmainfont{Latin Modern Roman}[Ligatures=TeX]
\setsansfont{Latin Modern Sans}
% Fonts bundled with Windows; no separate Amiri/Noto installation required.
\newfontfamily\arabicfont{Arial}[Script=Arabic,Scale=1.2]
\newfontfamily\hindifont{Nirmala UI}[Script=Devanagari,Scale=1.1]
\setmainlanguage{english}
\setotherlanguages{arabic,hindi}

% XeTeX primitives: change color at digit boundaries without opening groups.
\newXeTeXintercharclass\myDigitsClass
\XeTeXcharclass `0 = \myDigitsClass
\XeTeXcharclass `1 = \myDigitsClass
\XeTeXcharclass `2 = \myDigitsClass
\XeTeXcharclass `3 = \myDigitsClass
\XeTeXcharclass `4 = \myDigitsClass
\XeTeXcharclass `5 = \myDigitsClass
\XeTeXcharclass `6 = \myDigitsClass
\XeTeXcharclass `7 = \myDigitsClass
\XeTeXcharclass `8 = \myDigitsClass
\XeTeXcharclass `9 = \myDigitsClass
\XeTeXinterchartoks 0 \myDigitsClass = {\color{red}}
\XeTeXinterchartoks \myDigitsClass 0 = {\color{black}}
\XeTeXinterchartoks 4095 \myDigitsClass = {\color{red}}
\XeTeXinterchartoks \myDigitsClass 4095 = {\color{black}}

\title{\textbf{XeLaTeX Advanced Document}}
\author{Automated Production Framework}
\date{October 2026}
\begin{document}
\maketitle
\section{XeTeX Engine Test}
\begingroup
\XeTeXinterchartokenstate=1
Native XeTeX digit coloring: 2026, 847, and 100000.
\par
\endgroup
\section{Multilingual Script Rendering}
\begin{itemize}
  \item \textbf{Right-to-left Arabic}
    \begin{Arabic}
      الجمهورية والخطوط العربية المعقدة يتم تنسيقها تلقائياً وبشكل صحيح.
    \end{Arabic}
  \item \textbf{Devanagari ligatures}
    \begin{hindi}
      देवनागरी लिपि और जटिल संयुक्ताक्षर (जैसे: श्री, ज्ञान, क्र) रेंडर होते हैं।
    \end{hindi}
\end{itemize}
\section{Geometric Drawing}
\begin{center}
  \begin{tikzpicture}[scale=0.8]
    \draw[thick,color=blue!60,fill=blue!5] (0,0) circle (2cm);
    \node[align=center] at (0,0) {Native\\Unicode\\Shaping};
    \draw[line width=2pt,->,color=orange] (2.5,0) -- (4.5,0);
    \node[draw,rectangle,rounded corners,fill=green!10] at (6.5,0) {\textbf{PDF Generation}};
  \end{tikzpicture}
\end{center}
\end{document}
